\section{ベイズ的判断の方法}
\label{b27_bayes_judgement}

\subsection{授業内容}

\subsubsection{概要}
ベイズ的な判断方法は，従来の頻度主義的な$p$値に基づく仮説検定とは異なる視点からデータを解釈し，研究者がより柔軟かつ情報量豊富な判断を下すための手段を提供する。本コマでは，まず帰無仮説検定のロジックを振り返り，ベイズ的判断との対比を確認したうえで，ベイズ推定の結果をもとにした具体的な判断方法を学ぶ。すなわち，確信区間とROPEに基づく判断，ベイズファクターを用いた仮説検証，そして事後予測分布を使った判断である。これらの方法を一通り見渡しつつ，それぞれが何を判断材料にしているのかを整理する。

\subsubsection{コマ主題細目}
\begin{description}
	\item[ベイズ法と帰無仮説検定] あらためて帰無仮説検定のロジックを振り返りつつ，ベイズ推定の考え方との相違について確認する。$p$値が「帰無仮説のもとでデータが得られる確率」であって「帰無仮説が正しい確率」ではないという基本的な誤解を整理し，ベイズ的アプローチでは「パラメータがどこにあるか」を直接に確率で語れるという根本的な違いを理解する。特に事前分布を活用した推定方法は，時として何の前提も置かない帰無仮説検定よりも有用であることに言及する。

	\item[確信区間とROPEによる判断] ベイズ推定では，パラメータの事後分布から確信区間（Credible Interval）を算出できることを学ぶ。例えば95\%確信区間は，「パラメータが区間内に95\%の確率で存在する」と直接的に解釈できる。これは頻度主義的な信頼区間とは解釈が異なり，より直感的である。また，実質的等価性の領域（ROPE: Region of Practical Equivalence）の概念を導入し，パラメータがゼロや特定の値と「実質的に等価」かどうかを判断する方法を学ぶ。具体的には，95\%確信区間がROPEに完全に含まれる場合は「実質的に等価」，完全に外れる場合は「実質的に異なる」，部分的に重なる場合は「判断保留」と判断できる。これにより，単なる「有意/非有意」の二分法ではなく，パラメータの大きさと不確実性を考慮した判断が可能になることを理解する。

	\item[ベイズファクターによる判断] ベイズ統計学における仮説検証の方法として，ベイズファクター（Bayes Factor, BF）の概念と解釈について学ぶ。ベイズファクターは，二つのモデルの相対的な証拠の強さを表す比率であり，$BF_{10} = \frac{p(data|H_1)}{p(data|H_0)}$と定義される。ここで$p(data|H)$は，モデル$H$のもとでデータが観測される確率（周辺尤度）である。$BF_{10} > 1$のとき対立仮説$H_1$が，$BF_{10} < 1$のとき帰無仮説$H_0$が支持される。$BF_{10}$の値の大きさは証拠の強さを表し，例えば$BF_{10} = 10$は「対立仮説のための証拠が帰無仮説のための証拠の10倍ある」ことを意味する。ベイズファクターの解釈基準（例：1-3は「わずかな証拠」，3-10は「中程度の証拠」，10以上は「強力な証拠」）について学び，具体的な研究例での解釈方法を理解する。実際の計算と実践は次コマのJASPに譲る。
	      \begin{flushright}
		      $\to$ \textcite{Lee201709}のP.88--103.
	      \end{flushright}

	\item[事後予測分布による判断] 確信区間・ROPEやベイズファクターのほかに，事後予測分布を使った判断という選択肢もある。事後分布から得られた予測データの分布を眺めることで，「ある群が別の群より高くなる確率」（優越率）や「観測値がどの程度の範囲に収まるか」（被覆率）といった，母集団パラメータではなくデータレベルでの判断ができるようになる。これらは確率的プログラミング言語を使うことで初めて自由に設計できる方法であり，詳細な実装は最終コマで扱う。本コマではあくまで「こうした判断方法もある」という見取り図にとどめる。

\end{description}
\subsubsection{キーワード}
\begin{itemize}
	\item 帰無仮説検定との比較
	\item 確信区間とROPE
	\item ベイズファクター
	\item 事後予測分布
\end{itemize}
\subsection{授業情報}
\paragraph{コマの展開方法}
講義
\paragraph{標準シラバスにおける位置づけ}
科目番号5；心理学統計法；(1)心理学で用いられる統計手法；J より高度な記述や推定を目指して
\subsubsection{予習・復習課題}
\paragraph{予習}
前回までに学んだベイズ推定の基礎（事後分布，確信区間，HDIなど）を復習しておくこと。特に以下の点について確認しておくとよい：
\begin{itemize}
  \item 事前分布と事後分布の関係，およびベイズの定理による更新の仕組み
  \item 事後分布の要約方法（EAP, MAP, MED, HDI）の意味と解釈
  \item 頻度主義的な信頼区間とベイズ的な確信区間の解釈の違い
\end{itemize}
また，帰無仮説検定における$p$値の解釈の問題点（例：$p$値は「帰無仮説が正しい確率」ではない）について振り返っておくと，ベイズ的判断の利点がより理解しやすくなる。参考文献として\textcite{Lee201709}のP.88--103を事前に目を通しておくことが望ましい。

\paragraph{復習}
授業で学んだ内容を定着させるために，以下の点について理解を深めておくこと：
\begin{enumerate}
  \item 確信区間とROPEに基づく判断方法について，「実質的に等価」「実質的に異なる」「判断保留」の3つの判断がどのような場合に下されるか，具体例を挙げて説明できるようにする
  \item ベイズファクター$BF_{10}$の定義と解釈について，$BF_{10}=3$や$BF_{10}=0.1$の場合にそれぞれ何を意味するか説明できるようにする
  \item ベイズファクターの解釈基準（1--3: わずかな証拠，3--10: 中程度，10以上: 強力な証拠）を覚え，研究報告での使い方を理解する
  \item 頻度主義的な$p$値による判断と，ベイズ的判断（確信区間/ROPE/ベイズファクター/事後予測分布）の違いを，自分の言葉で説明できるようにする
\end{enumerate}
次回のJASPを用いた実習に備えて，JASPを各自のPCにインストールしておくこと（https://jasp-stats.org/）。
